\honest{Where Recursive Justification Hits Bottom}{Eliezer Yudkowsky}{2008}

Why do I believe the Sun will rise tomorrow? Because it has risen on thousands of days before. But why do I believe the future will be like the past? The laws of physics and Occam's razor do not help, since they too have only worked so far. \nb{The sunrise is Hume's example, in section 4 of Hume's \textsc{Enquiry}. The notes on ``A Priori'' set out Hume's argument that any defense of induction assumes it.}

Some people conclude that science rests on faith too, ``so don't you criticize me.'' I answered this before, in ``The Fallacy of Gray,'' by asking why people who praise faith say it so triumphantly. But I admit that this does not answer the real problem: if every belief needs a justification, how does the chain end? And if it may end in an assumption, why not assume anything? \nb{The classic answer, foundationalism, holds that some beliefs are justified without resting on other beliefs, so the chain ends in something justified, not something assumed. For induction that would mean an a priori justification, which Yudkowsky had argued against in ``A Priori.'' The post does not take up this option.}

Bayesian updating does not solve it either, I admit. After 3 red and 6 white balls, the chance that the next ball is red is 4/11 on one prior and 7/11 on another. There are possible minds with anti-Occamian and anti-Laplacian priors, and when asked why they keep using them, they say: ``Because it's never worked for us before!'' \nb{Both figures are right. The anti-Laplacian mind is the standard objection to justifying induction by induction: counterinduction can justify itself the same way.}

My approach is to treat ``Should I trust my brain?'' and ``Should I trust Occam's Razor?'' as ordinary questions. My brain came from natural selection, which gives reasons to doubt it, but would also have removed brains with priors that bad. I go on examining my rules, with my current brain, because there is nothing else to use. \nb{Quine made the same evolutionary argument in 1969: creatures ``inveterately wrong in their inductions'' tend to die out before they reproduce.}

So when someone keeps asking why I believe what I believe, I go around in a loop: induction has usually worked, simplicity explains why, and evolution explains how my brain could see this. Have I licensed circular logic? No, I say; I have licensed reflecting on my mind with my mind. A sheet of paper that says ``Everything on this sheet of paper is true'' gives no causal story of how it could match the world. \nb{Philosophers call the post's kind of loop ``rule-circular'': it uses a rule to show that the rule is reliable. The sheet of paper is ``premise-circular'': its conclusion is among its premises. Some philosophers hold that only the second kind is vicious. The distinction supports the post, which does not use these terms.}

Would most minds that settle into coherence be incorrect? Perhaps, but we were not generated at random; we evolved.

Is this the believer who trusts the Bible because the Bible says so? People who leave religion do not jump to a state of pure emptiness. They notice unanswered prayers and a God who cannot tell them the hundredth digit of pi, and they question their beliefs with the minds they have. \nb{This is Otto Neurath's image of 1932: sailors who must rebuild their ship on the open sea, with no dry dock in which to take it apart.} If the Bible had been right about everything else, its claim to come from God would be worth taking seriously. What condemns it is its record, not its loop.

What is the alternative to induction, I ask: to believe that the future will not be like the past, because that has always failed before? Rationalists are not out to win arguments with ``ideal philosophers of perfect emptiness''; we are ``out to win.'' \nb{Hume went on reasoning from experience too, but called it custom, not reasoning. The post agrees with Hume in practice and differs in counting its loop as a justification.}

``Everything, without exception, needs justification.'' Some justifications loop, and ``I do think'' loops can be told from circular logic ``by common sense.'' \nb{The post suggests tests along the way, a causal story of how the mind tracks the world and a record on other matters, but states no criterion. Two days later, in ``My Kind of Reflection,'' Yudkowsky wrote: ``I haven't yet sat down and formalized the exact difference.'' Yudkowsky also wrote there of being ``far from the first person to consider reflective application of reasoning principles.'' Nelson Goodman had argued in 1955 that rules of induction are justified by agreement with the inferences we accept.} Loops grant no ``immunity from questioning.'' Apply full force, loop or not, and play to win. \nb{W. W. Bartley gave a similar reply to the argument that reason, like faith, rests on an arbitrary commitment: a view that holds everything open to criticism can itself be held open to criticism. Unlike the post, Bartley dropped the demand that beliefs be justified.}
